\subsection{Beantwortung der Fragestellungen}
\label{subsec:beantwortung-fragestellungen}

\paragraph{Fragestellung 1: Adressierte Anforderungen.}
Das Template adressiert die Anforderungen durch eine generierbare Baseline,
klare Komponenten, Profile, gemeinsame Workflows, zentrale Pflege und
versionierte Konfiguration. Seit Governance~v1 werden ausgewählte Größen-,
Komplexitäts- und Abhängigkeitsregeln zusätzlich über zentrale Policy, lokale
CLI und CI operationalisiert. Getrennte Laufzeit- und Prüfpfade stützen dies.
Fachlogik, Authentifizierung, Berechtigungen und Produktivbetrieb bleiben
produktspezifisch; eine vollständige semantische Architekturprüfung bleibt
Reviewaufgabe.

\paragraph{Fragestellung 2: Varianten für Web, Cloud und Desktop.}
Fünf Presets kombinieren das gemeinsame Frontend mit Backend-, Cloud- und
Tauri-Bausteinen; PostgreSQL wird nur bei Backendfähigkeit ergänzt. Generierte
Projekte enthalten überwiegend benötigte Pfade. \texttt{desktop-cloud} und
\texttt{full-platform} sind technisch gleich generiert und semantisch
abgegrenzt.

\paragraph{Fragestellung 3: Wiederverwendung.}
Die zentrale Pflege von Implementierungen, Profilen, Generator, Konfiguration,
Tests und Dokumentation fördert Wiederverwendung und einheitliche
Ausgangsstände. Der Vorteil gilt für die Baseline und künftige Generierungen;
bestehende Produktprojekte erhalten keine automatischen Updates.

\paragraph{Fragestellung 4: Reduktion wiederkehrender Aufwände.}
Automatisierte oder einheitliche Profilwahl, Generierung, Diagnose,
Installation, Start, Quality, Tests und Builds schaffen plausibles
Einsparpotenzial. Rule IDs und Handlungshinweise vereinheitlichen die
Rückmeldung für Menschen und Coding Agents. Zeit-, Aufwands- und
Vergleichsdaten fehlen jedoch; eine Produktivitätssteigerung oder geringere
Fehlerquote ist unter Einbezug zentraler Pflege und mehrerer Toolchains nicht
nachgewiesen.

\paragraph{Fragestellung 5: Projekteignung.}
Das Template eignet sich als Basis für verwandte Web-, Desktop- und
kombinierte Business-Anwendungen mit abbildbarem Plattform- und Cloudbedarf.
Persistenzbedarf, Persistenzprovider und Source of Truth werden davon getrennt
produktspezifisch bestimmt. Stark native Anwendungen, Spiele und spezialisierte
Echtzeitsysteme liegen außerhalb des primären Zielbereichs. Die qualitative
Einstufung ist keine universelle Eignungszusage.

\paragraph{Fragestellung 6: Grenzen, Wartungsaufwände und Risiken.}
Grenzen sind wachsende Regressionsbreite, manuelle Vertragssynchronisierung,
die nicht strikt getrennte Capability-Auswahl, fehlende Produktupdates und der
Pflegebedarf mehrerer Toolchains. Für die Governance kommen technisch
unterdrückbare \texttt{CQ001}-Befunde, unvollständig validierte Exceptions,
heuristische Architekturregeln, fehlende negative Rust-Fixtures und die
dateibasierte statt manifestbasierte Profilableitung hinzu. Backendhaltige
generierte Profile, der lokale Rust-/Tauri-Pfad und Windows sind nicht grün;
Release Validation, Compose-Start, Dokumentationssynchronität und Branch
Protection fehlen oder sind offen.
Signing, Notarisierung, Updater, Produktiv-Deployment und Sicherheitsmodell
bleiben Produktaufgaben.
